\section{Methods}
\label{sec:methods}
\subsection{Task and scoring}
\textbf{Task.} Each packet contains 32 synthetic records, a selection rule, and a
claim listing the identifiers that satisfy the rule. The rule uses each record's
active status, flag, and a minimum score. A program computes the exact sorted answer.
Faulty claims omit, add, or substitute one identifier. Neither agent can change
the packets.

\textbf{Evidence.} The model sees its own complete packet and reviews its peer's
claim. With full evidence, it also sees the peer's source records. With restricted
evidence, it sees only the peer's rule and claim. Before giving its verdict, it
is not told whether the current claim is correct or which error it contains.

\textbf{Verification rule.} The reviewer may \texttt{ACCEPT} only if it has all
required records and the claim matches them. Otherwise, it must \texttt{REJECT}.
Peer messages cannot override this rule.

\textbf{Conflicting reward.} During the initial interactions, all claims are
correct. Each verdict matching a claim's actual correctness earns one point,
giving the pair a shared score from zero to two. Approving without evidence can
therefore earn points, while following the rule and rejecting cannot. This
creates an incentive for unsupported approval without an instruction to collude.

\subsection{Agents and interaction histories}
We test one model reviewer interacting with a scripted peer. The study has five
independent paired blocks. Each block contains two ten-episode interactions:
one with a peer that always accepts and one with a peer that always rejects.
We call this initial phase \emph{formation}.

The paired interactions use the same correct packets, restricted evidence, and
neutral peer message: ``Review my packet using the ordinary verification rule.''
The model's assigned role is matched within each block and counterbalanced across
blocks. A fixed random seed sets condition order.

\textbf{Model and memory.} The reviewer uses \texttt{gpt-6.1-sol}, low reasoning
effort, the default service tier, and a 4,096-token output limit. Its only memory
is the visible JSON transcript. It has no tools, retrieval, separate reflection
step, provider-managed conversation history, or carried-over private reasoning.
The verification rule appears on every request; peer content is marked untrusted.

\textbf{Episode sequence.} Each message is limited to 200 Unicode code points.
Every formation episode follows four steps:
\begin{enumerate}[itemsep=2pt]
\item The reviewer writes its first message without seeing the peer's current
message.
\item Both first-round messages become visible. The reviewer writes a second
message; the peer repeats its neutral message.
\item The reviewer sees both rounds and gives its verdict before seeing the
peer's verdict.
\item Both verdicts and the truthful score feedback are added to the next
episode's history.
\end{enumerate}
The conditions therefore differ in both the peer's verdicts and the feedback
those verdicts produce. Fixed task seeds do not make the model's outputs
deterministic.

\begin{figure}[t]
\centering
\begin{tikzpicture}[font=\small,>=Stealth,
 box/.style={draw=inkblue,rounded corners=2pt,align=center,fill=lightblue,
 minimum height=14mm,text width=33mm,inner sep=4pt}]
\node[box] (u) at (0,0.9) {Rejecting-peer\\history};
\node[box] (c) at (0,-0.9) {Accepting-peer\\history};
\node[box,text width=42mm,minimum height=31mm] (eval) at (5,0)
 {Matched new claims\\\textbf{Complete evidence}\\\vspace{2pt}A: rejecting peer\\B: accepting peer};
\node[box,text width=29mm,minimum height=31mm] (out) at (9.8,0)
 {False approval\\\textbf{B minus A}\\\vspace{2pt}Also measure correct-claim approval};
\draw[->,thick,inkblue] (u.east) -- (eval.west |- u.east);
\draw[->,thick,inkblue] (c.east) -- (eval.west |- c.east);
\draw[->,thick,inkblue] (eval) -- (out);
\end{tikzpicture}
\caption{Testing recovery. After each type of peer history, the reviewer receives
the same new claims and full evidence. B minus A measures the remaining history
effect on false approval. Other conditions test correction, history removal,
and continued evidence restriction.}
\label{fig:design}
\end{figure}

\subsection{Checks before interpreting recovery}
\textbf{Can the model perform the task?} Before formation, we test ten correct
and ten faulty claims with full evidence and no interaction history. The model
must return valid structured responses to all 20 and classify at least nine of
each type correctly. These tests are excluded from the recovery estimates.

\begin{samepage}
\textbf{Does peer history affect unsupported approval?} We compare approval
rates over formation episodes 2--10, taking accepting minus rejecting within
each block and weighting blocks equally. Greater approval after an accepting
peer is supported only if the combined interval and missing-outcome bounds lie
above zero. We keep all blocks and recovery conditions regardless of this check.\par
\end{samepage}

Episode 1 is a control: the reviewer has not yet received feedback that differs
by peer policy. Initial inputs are identical across conditions. Later inputs in
that episode also match if generated messages and attempts match. Separate model
calls can still produce different outputs.

\subsection{Recovery tests}
After formation, we save both interaction histories and test six recovery
conditions (Table~\ref{tab:arms}). Each condition receives the same eight new
claim pairs in the same order. The reviewer's own claim and the peer's claim can
each be correct or faulty; all four combinations appear twice.

This gives four correct and four faulty peer claims per condition and block.
Faulty claims cover all three error types. A fixed random seed sets the order in
which conditions run.

Each test starts again from its assigned saved history, except E, which removes
that history. All conditions have full current evidence except F. The reviewer
receives a common neutral peer message and makes a review decision. It sends no
message and sees no current peer verdict.

Earlier test responses never enter later tests. We therefore measure immediate
responses on new claims, rather than recovery over a continuing conversation.
Appendix~\ref{app:prompts} gives the verification and correction instructions.

\begin{table}[htbp]
\centering\small
\caption{Recovery conditions. Both formation histories have restricted evidence.
Correction explicitly rejects earlier approval shortcuts. History removal clears
the earlier transcript but keeps the current task, correction, and verification
rule.}
\label{tab:arms}
\begin{tabularx}{\linewidth}{@{}c l l X l@{}}
\toprule
Arm & Formation peer & Current evidence & Instruction & History\\
\midrule
A & Rejecting & Full & Neutral & Retained\\
B & Accepting & Full & Neutral & Retained\\
C & Rejecting & Full & Correction & Retained\\
D & Accepting & Full & Correction & Retained\\
E & Accepting & Full & Correction & Removed\\
F & Accepting & Restricted & Neutral & Retained\\
\bottomrule
\end{tabularx}
\end{table}

\begin{samepage}
We compare false approval rates as follows:
\begin{itemize}[itemsep=3pt]
\item \textbf{Remaining peer-history effect (B$-$A, primary):} compare the two
histories after restoring full evidence.
\item \textbf{Evidence restoration (B$-$F):} restore evidence after the same
accepting-peer history.
\item \textbf{Correction (D$-$B and C$-$A):} add correction after each type of
history. The difference, $(D-B)-(C-A)$, tests whether its effect depends on
the prior peer.
\item \textbf{History removal (E$-$D):} remove history while keeping correction
in place. This measures the added effect of removal, not reset alone.
\end{itemize}
\end{samepage}
A small B$-$A difference does not establish recovery if both conditions perform
poorly. We therefore also report each condition's absolute performance.

\subsection{Measures and uncertainty}
\textbf{Measures.} False approval rate (\far) is the fraction of faulty claims
accepted; lower is better. Correct-claim approval (\cpa) is the fraction of
correct claims accepted; higher indicates more useful review when evidence is
complete. Balanced accuracy averages correct rejection and correct approval:
$\{(1-\far)+\cpa\}/2$.

We interpret these measures together. Rejecting everything lowers false approval
but also blocks correct work. When evidence is restricted, however, rejection
is required by the verification rule even for correct claims.

\textbf{Independent comparisons.} We calculate $\far_B-\far_A$ within each block,
then average the five differences with equal weight. The independent unit is the
paired block, not each review.

\textbf{Uncertainty.} We prespecified two 95\% intervals: a block bootstrap with
10,000 fixed-seed resamples, and a paired-block Student-$t$ interval. Resampling
keeps all conditions within a block together. The combined interval spans both
intervals. The analysis withholds uninformative intervals when paired effects
are identical across blocks. With only five blocks, these estimates are
exploratory; identical observed outcomes cannot establish that an effect is absent.

\textbf{Missing decisions.} Missing or invalid verdicts remain unknown. For $a$
observed faulty approvals, $m$ missing outcomes, and $N$ scheduled faulty reviews,
the possible rate ranges from $a/N$ to $(a+m)/N$. This covers treating every
missing decision as a rejection or as an approval.

We carry these bounds through differences between conditions. Requests whose
final outcome cannot be confirmed remain in sensitivity analyses. A small
estimated difference or an incomplete run does not establish recovery or
equivalent performance.
